\documentclass[11pt]{article}
\usepackage[a4paper,margin=25mm]{geometry}
\usepackage{amsmath,amssymb,amsthm,mathtools,booktabs,hyperref,xurl}
\setlength{\emergencystretch}{2em}
\newtheorem{theorem}{Theorem}
\newtheorem{lemma}[theorem]{Lemma}
\newtheorem{proposition}[theorem]{Proposition}
\newcommand{\Z}{\mathbb Z}\newcommand{\Q}{\mathbb Q}
\title{An odd second-Pell quotient removes one inverse}
\author{\(\mathbb Z\mathbb Q\) Diophantine Project}
\date{24 November 2026}
\begin{document}
\maketitle
\begin{abstract}
The standard positive completion of the second Pell equation already
permits its nonzero integer quotient to be odd. Certifying its half-shift
therefore makes nonvanishing automatic and removes one reciprocal
coordinate. Composing this observation with the common-numerator
ordinary-\(abc\) table gives 674 rational unknowns and 593 quadratic
residuals. Every retained arithmetic premise and the external Jones
coding theorem remain required. The equivalence is on source-instance
projections, since the integer-certificate witnesses must be chosen for
a different integer input.
\end{abstract}

\section{Oddness in the existing positive completion}

For integers \(a\ge2\), define
\[
 X_n(a)+Y_n(a)\sqrt{a^2-1}
       =(a+\sqrt{a^2-1})^n,\qquad n\ge0.
\]
Both sequences satisfy \(Z_{n+2}=2aZ_{n+1}-Z_n\), with
\(Y_0=0,Y_1=1\). The binomial certificate of~\cite{Binomial}
and the one-quotient third-Pell completion of~\cite{CRT} use
\(B=2r+1\ge3\) and
\[
              c=Y_B(a),\qquad d=X_B(a).
\]
The second Pell norm is
\begin{equation}\label{norm}
                F^2-(a^2-1)i^2c^4=1.
\end{equation}

\begin{lemma}\label{odd}
For every integer \(a\ge2\) and odd integer \(B\ge3\), put
\(c=Y_B(a)\), \(d=X_B(a)\), and \(q=Bc\). Then
\[
 F=X_q(a),\qquad i=Y_q(a)/c^2
\]
are positive integers satisfying \eqref{norm}, with \(i>1\)
and \(i\) odd.
\end{lemma}
\begin{proof}
Modulo two the recurrence for \(Y_n\) reads
\(Y_{n+2}\equiv Y_n\), so \(Y_n(a)\equiv n\pmod2\).
Consequently \(c,q,Y_q(a)\) are odd.

Put \(\Delta=a^2-1\). Expanding
\[
 (a+\sqrt\Delta)^{Bc}=(d+c\sqrt\Delta)^c
\]
and comparing second coordinates gives
\[
 Y_q(a)=\sum_{\substack{1\le j\le c\\j\ {
\rm odd}}}
       \binom cj d^{c-j}c^j\Delta^{(j-1)/2}.
\]
The \(j=1\) summand is \(c^2d^{c-1}\), and every other
summand contains \(c^3\). Thus \(c^2\mid Y_q(a)\).
Since both numerator and divisor are odd, \(i\) is odd.
Moreover all summands are positive, giving
\(i\ge d^{c-1}>1\): indeed \(c=Y_B(a)>B\ge3\) and
\(d=X_B(a)>1\). The Pell norm for index \(q\) gives
\eqref{norm}.
\end{proof}

The index \(q=Bc\) is already the explicit choice in the positive
completion of~\cite{Binomial}. Thus oddness introduces no new
divisibility or padding hypothesis. The positive third-Pell completion
in~\cite[Proposition 3]{CRT} applies to every positive \(i,F\)
satisfying \eqref{norm}, so it applies to this choice. Its third index
may still be increased in the prescribed arithmetic progression to make
the remaining quotient positive.

\section{The new integer input and the one-row replacement}

The complete predecessor is the 675-coordinate, 594-row
ordinary-\(abc\) table of~\cite{Common}. It has a certified
integer input \(i\), a private rational coordinate \(v\), and
two rows
\begin{equation}\label{old}
                         iv=1,\qquad vE=C_2.
\end{equation}
Its retained integer-input chart for \(c\) supplies \(C_2=c^2\).
The coordinate \(v\) occurs nowhere else.

Replace the certified integer input \(i\) by an integer input \(h_i\)
and interpret the actual source quotient as
\begin{equation}\label{new-i}
                              i=2h_i+1.
\end{equation}
Delete \(v\) and both rows in \eqref{old}, and insert
\begin{equation}\label{new}
                           E=(2h_i+1)C_2.
\end{equation}
This is a quadratic polynomial equation in the retained coordinates.
All other source rows are unchanged. The former \(i\)-certificate
now certifies \(h_i\), with the same shared elliptic coordinates and
the same contact offset. Its internal role labels are retained to make
the table comparison transparent.

\begin{theorem}\label{projection}
Under exactly the arithmetic-interface premises of the predecessor,
the replacement \eqref{new-i}--\eqref{new} preserves its source-instance
projection. The new complete table has 674 rational unknowns and 593
quadratic residuals, with all thirteen integer predicates retained.
Its sum of squared residuals has degree exactly four, and has the same
conditional rational-undecidability interpretation as the predecessor.
\end{theorem}
\begin{proof}
For soundness, the retained integer interface makes \(h_i\in\Z\).
Therefore \(i=2h_i+1\) is an odd, nonzero integer. The retained
\(c\)-square row gives \(C_2=c^2\), and \eqref{new} gives
\(E=ic^2\). Thus the old outer relations \eqref{old} hold after
setting \(v=1/i\). In particular the second Pell norm and every
other integer-source hypothesis are recovered in the same order as
in the preceding signed-source proof. That proof yields the required
coding instance. This argument uses the soundness of the integer
interface directly; it does not assert that an elliptic certificate
for \(h_i\) is a certificate for \(i\).

For completeness, start from any member of the source's represented
set. Choose the positive corrected binomial witnesses as before,
including \(c=Y_B(a)\). Lemma~\ref{odd} gives a positive odd
\(i>1\), so \(h_i=(i-1)/2\) is a positive integer.
The positive third-Pell completion supplies the remaining quotient,
and the other positive source coordinates and the strengthened
margin are obtained as in the predecessor. The choices of \(i,F\)
do not enter that margin or the short fixed-power certificate.

Apply the retained simultaneous integer-interface completeness theorem
to the new thirteen-input tuple, using \(h_i\) in place of \(i\).
Every input is nonzero; this verifies the required completeness
hypothesis. Choose its common numerator and all auxiliary elliptic
witnesses afresh. Equation \eqref{new} now holds with \(E=ic^2\),
and every other source row holds. This proves equality of instance
projections, without claiming that an arbitrary old witness has an
odd quotient or that its old elliptic coordinates survive unchanged.

One coordinate is deleted. Two rows are replaced by one, and the
coordinate relabeling adds no coordinate. Hence the new counts are
\(675-1=674\) and \(594-1=593\). All rows have degree at most
two, and at least one unchanged row has nonzero quadratic homogeneous
part. Over \(\Q\), a sum of squares vanishes precisely when every
summand vanishes; over \(\mathbb R\), the sum of squares of the
nonzero leading quadratic parts is not identically zero. Thus the sum
has degree exactly four. The source and arithmetic premises are
unchanged.
\end{proof}

\section{Exact implementation and verification scope}

The deleted coordinate is \texttt{i\_inverse}; the old coordinate
\texttt{i} is relabeled \texttt{i\_half}. The two old rows are
\texttt{integer.i\_reciprocal.unit} and
\texttt{jones1980.E.reciprocal}. The new row is
\[
 \mathtt{E}-2\,\mathtt{i\_half}\cdot\mathtt{int\_c\_29}
                     -\mathtt{int\_c\_29}=0.
\]
The full occurrence audit finds that the old \texttt{i} occurs only
in its own square row and the deleted unit row. The inverse occurs
only in the two replaced rows. Every other row agrees exactly after
the coordinate relabeling. The surviving \texttt{integer.i.*} labels
and \texttt{int\_i\_reciprocal\_*} names are historical interface
role names, not assertions that a reciprocal remains.

The independent verifier imports no builder. It checks the complete
row comparison, coefficient ring, degree, support, variable counts,
input relabeling and source qualifications. It also checks six exact
Pell fixtures by separate powering and recurrence calculations. These
fixtures illustrate the quantified proof of Lemma~\ref{odd}; they
are not full witnesses for the universal coding prefix.

The deterministic scripts, full coefficient table and certificates
are in \nolinkurl{undecidability/odd_pell_quotient_v104/}.
Run from the repository root:
\begin{verbatim}
python -B undecidability/odd_pell_quotient_v104/build_odd_quotient.py
python -B undecidability/odd_pell_quotient_v104/verify_odd_quotient.py
\end{verbatim}

No general or specified-interface numerical saving is asserted here:
those routes use genuine reciprocal integrality charts, and their
coordinate and nonzero bookkeeping differ from the ordinary-input
\(abc\) chart used in this theorem.

\begin{thebibliography}{9}
\small
\bibitem{Binomial} Project, \emph{An exact Pell certificate for a power
of two and central-binomial divisibility}, 23 November 2026.\par
\nolinkurl{undecidability/papers/binomial_pell_certificate.tex}.
\bibitem{CRT} Project, \emph{A single integer quotient for the third
Pell certificate}, 23 November 2026.\par
\nolinkurl{undecidability/pell_crt_v95/one_quotient_pell.tex}.
\bibitem{Common} Project, \emph{A common numerator compresses reciprocal
products in the ordinary-abc quartic}, 24 November 2026.\par
\nolinkurl{undecidability/abc_common_numerator_v103/common_numerator.tex}.
\end{thebibliography}
\end{document}
